\begin{proof}
First suppose that $\mathcal X$ is split, and fix a point of the ordinary
open subspace under consideration.  Choose an affine toric chart containing
its $\mathcal X$-coordinate, and let $\tau$ be the face corresponding to
the orbit of that coordinate.  The principal-open decomposition
\eqref{eq:p061-orbit-product-chart} identifies a neighbourhood of the orbit
with
\[
 X_{\tau,N_\tau}^{\mathrm{an}}\times_KT_{N'}^{\mathrm{an}},
\]
where $\tau$ is full-dimensional and strongly convex in
$(N_\tau)_{\mathbf R}$, and the chosen point has first coordinate the
torus-fixed point of $X_{\tau,N_\tau}$.

Put
\[
 Y_\tau=T_{N'}^{\mathrm{an}}\times_KY.
\]
This is a good strict analytic space locally at the chosen point, and its
reduction is quasi-smooth.  Indeed, on strict affinoid charts,
\Cref{lem:p061-product-factor-reduction} applied to the zero ideal on the
torus chart and the nilradical of the $Y$-chart gives a nilpotent closed
immersion
\[
 T_{N'}^{\mathrm{an}}\times_KY_{\mathrm{red}}
 \longrightarrow Y_\tau.
\]
Its source is quasi-smooth, hence reduced, because a split torus and
$Y_{\mathrm{red}}$ are quasi-smooth and quasi-smoothness is stable under
finite products.  It is therefore the reduction of $Y_\tau$.

Apply \Cref{prop:p061-parameterized-toric-fixed-fibre} with parameter
\[
 T_{N'}^{\mathrm{an}}\times_K C_1\times_K\cdots
      \times_K C_m\times_KY.
\]
After shrinking to a product neighbourhood, the given function is the
pullback of a classically pluriharmonic function $g$ on that parameter.
Regrouping its factors as
\[
 C_1\times_K\cdots\times_KC_m\times_KY_\tau
\]
puts $g$ in the scope of
\Cref{thm:p061-arbitrary-curve-factors-ph-rpl}.  Thus $g$ is locally
$\mathbf R$-PL and Gubler--Jell--Rabinoff harmonic.  Pullback stability of
$\mathbf R$-PL functions
\cite[Lemma~5.4(1)]{GJR26} gives the same conclusion for the original
function, while \Cref{prop:ch18-harmonic-functorial} gives its harmonicity.
The argument is local at every point, and transfer through an analytic local
isomorphism proves the local-model assertion.

Now let $X$ become split over $L$.  The scalar-extension assertion used in
the proof of \Cref{cor:p061-algebraic-split-toric-forms} makes the pulled-back
function classically pluriharmonic.  The base-changed curve factors remain
good strict analytic curves, while the reduction of $Y_L$ is quasi-smooth;
hence the split case over $L$ makes the pulled-back function $\mathbf R$-PL.
Since $L$ is the completion of an algebraic extension, ground-field descent
\cite[Corollary~5.12]{GJR26} makes the original function $\mathbf R$-PL.
The final two-sign comparison in the proof of
\Cref{cor:p061-algebraic-split-toric-forms} then gives harmonicity over $K$.
Local analytic isomorphisms again preserve the conclusion.
\end{proof}
